% ===== 5 形貌识别与异常检测 =====
\section{形貌识别与异常检测}
\label{sec:shape-anomaly}

在完成主粒径分析链后，赛题还要求给出颗粒形状类别并识别明显异常物。本文不为这两项任务重新建立独立的视觉网络，而是尽可能复用第\ref{sec:geometry-features}节已经获得的颗粒级轮廓与几何特征。这样既减少现场模型数量，也使每一个形貌或异常判断都能回溯到具体实例及其测量量。需要特别说明的是，单目俯视图不能获得颗粒厚度，因此本章所称“针片状”仅表示\textbf{二维投影形貌候选}；同样，当前泥团判断主要是几何启发式筛查，不能等同于已经经过纹理与材质真值验证的语义分类器。

\subsection{投影形貌表征}
\label{sec:projected-morphology}

颗粒投影形貌可以从三个互补角度描述：伸长程度、边界圆滑程度以及轮廓相对于凸包的凹陷程度。本文直接复用式\eqref{eq:shape-descriptors}中的长宽比 $AR_i$、圆度 $C_i$ 和凸度 $S_i$。其中 $AR_i=b_i/a_i$ 越小，说明投影越细长；圆度
\begin{equation}
 C_i=\frac{4\pi A_i}{P_i^2}
 \label{eq:morph-circularity}
\end{equation}
在理想圆形时趋近于 1，并随着轮廓伸长或边界复杂度增加而降低；凸度
\begin{equation}
 S_i=\frac{A_i}{A_i^{\mathrm{convex}}}
 \label{eq:morph-solidity}
\end{equation}
则描述实际投影面积相对于凸包面积的比例，轮廓凹陷或棱角效应增强时通常会降低。

这三个指标的作用并不相同。仅依赖 $AR_i$ 可以筛出明显细长颗粒，却无法区分近方形与圆形颗粒；仅依赖 $C_i$ 又会同时对颗粒真实形状和分割边界粗糙度敏感；$S_i$ 可以补充描述凹陷与非凸性，但对于整体凸而有明显尖角的颗粒并不总是足够。因此，本文使用三者组合建立透明的规则基线，而不是将任一单指标直接解释为完整形貌真值。由于这些特征均来自实例掩膜，后续还可以通过人工专家标签重新标定阈值，而无需改变上游分割模型。

更重要的是，二维投影形貌与工程规范中的三维形状概念必须区分。真实针片状判定通常涉及颗粒不同空间方向尺寸，而单张俯视图只能观测两个投影主轴。因此，本文输出类别命名为“针片状候选”，其含义是 $AR_i$ 显著偏小的细长投影实例；最终若要报告规范意义下的针片状准确率，需要将图像结果与人工量规、三维测量或专家真值进行配对验证。

\subsection{可解释的形貌规则基线}
\label{sec:morphology-rules}

当前实现将颗粒投影形貌划分为“针片状候选、圆形、棱角状、普通”四类。设阈值集合为 $\boldsymbol{\tau}_{\mathrm{morph}}$，默认规则写为
\begin{align}
 \text{needle/flaky candidate:}\quad & AR_i<0.50, \label{eq:shape-needle}\\
 \text{round:}\quad & C_i>0.85,\quad S_i>0.95,\quad AR_i>0.80, \label{eq:shape-round}\\
 \text{angular:}\quad & (C_i<0.60\ \lor\ S_i<0.90)\quad\land\quad AR_i\ge0.50. \label{eq:shape-angular}
\end{align}
不满足上述条件的颗粒归入普通类。代码中使用固定的确定性判定顺序，并显式排除将细长颗粒同时归入棱角类，从而保证相同输入始终得到相同标签。表\ref{tab:morph-thresholds}给出了当前默认阈值及其作用。

\begin{table}[htbp]
\centering
\caption{当前投影形貌分类的默认阈值}
\label{tab:morph-thresholds}
\small
\renewcommand{\arraystretch}{1.15}
\begin{tabular}{>{\raggedright\arraybackslash}p{2.5cm}>{\centering\arraybackslash}p{2.0cm}>{\raggedright\arraybackslash}p{8.3cm}}
\toprule
参数 & 默认值 & 作用 \\
\midrule
$\tau_{AR}^{\mathrm{needle}}$ & 0.50 & 识别明显细长投影，输出“针片状候选” \\
$\tau_C^{\mathrm{round}}$ & 0.85 & 圆形类别要求较高圆度 \\
$\tau_S^{\mathrm{round}}$ & 0.95 & 圆形类别同时要求轮廓接近凸形 \\
$\tau_{AR}^{\mathrm{round}}$ & 0.80 & 避免将细长椭圆误判为圆形 \\
$\tau_C^{\mathrm{angular}}$ & 0.60 & 低圆度作为不规则/棱角候选条件之一 \\
$\tau_S^{\mathrm{angular}}$ & 0.90 & 低凸度作为棱角候选条件之一 \\
\bottomrule
\end{tabular}
\end{table}

这些数值是当前工程原型的经验默认值，而不是由本竞赛骨料专家标注拟合得到的最终阈值。因此，本报告不会仅凭规则输出比例声称形貌分类已经达到某一准确率。正式实验应抽取具有代表性的颗粒实例，由人工或专家给出独立投影形貌标签，再用混淆矩阵和 Macro-F1 等指标评价四类结果；若阈值需要调整，则必须只在开发/校准子集上确定，并在独立测试子集上报告最终性能。

采用规则法作为当前基线的主要优势是可解释和易校准。若后续人工标注表明某些类别在 $(AR,C,S)$ 空间高度重叠，再引入基于颗粒 crop 的轻量分类器会更合理；反之，如果简单阈值已经能够稳定区分比赛关注的典型形貌，则无需为了模型复杂度本身增加现场训练和推理环节。形貌任务由此保持与第\ref{sec:design-principles}节“校准显式、部署可控”的总体设计原则一致。

\subsection{尺寸异常与疑似泥团筛查}
\label{sec:anomaly-detection}

异常检测需要区分两类性质不同的问题。第一类是具有明确物理阈值的\emph{尺寸异常}，例如赛题关注的 $>50$~mm 大块；这类判断可以直接建立在经过尺度恢复和筛分粒径映射后的 $d_i$ 上。第二类是泥团、杂物等\emph{材质或语义异常}，它们与正常骨料的差异可能主要体现在颜色、纹理、破碎结构或材质，而不一定能够由尺寸和轮廓唯一确定。本文当前实现对前者采用明确物理阈值，对后者只提供保守的几何候选规则，并在输出中使用“疑似”而非确定类别。

当前异常模块的尺寸规则为
\begin{align}
 d_i>50~\mathrm{mm} &\quad\Rightarrow\quad \text{大块异物/超限}, \label{eq:anomaly-oversize}\\
 d_i<5~\mathrm{mm} &\quad\Rightarrow\quad \text{过小实例/噪声}. \label{eq:anomaly-undersize}
\end{align}
其中 $<5$~mm 类主要服务于当前粗骨料分析口径下的实例清理和统计核查，并不意味着物理上所有小于 5~mm 的颗粒都属于“错误目标”；在包含细集料的应用场景中，该下限必须根据任务重新定义。类似地，$>50$~mm 判断的可靠性依赖第\ref{sec:scale-calibration}节的尺度标定和第\ref{sec:sieve-equivalent-diameter}节的筛分等效粒径映射，因此尺寸异常仍会继承上游尺度与分割误差。

对于疑似泥团，当前代码在 $40<d_i\le50$~mm 范围内使用凸度作为一个几何兜底条件：
\begin{equation}
 40<d_i\le50~\mathrm{mm}\quad\land\quad S_i<0.85
 \quad\Rightarrow\quad \text{疑似泥团/非骨料}.
 \label{eq:anomaly-mud}
\end{equation}
该规则的动机是将“大尺寸且轮廓明显不规则”的实例作为优先复核对象，但它并不能识别所有泥团，也可能将正常的不规则骨料误报为异常。实现允许通过 \texttt{mud\_solidity=None} 完全关闭这一兜底规则，说明该模块本身被设计为可选筛查器，而不是高置信度语义识别器。

\begin{table}[htbp]
\centering
\caption{当前异常检测类别、判据及证据强度}
\label{tab:anomaly-rules}
\small
\renewcommand{\arraystretch}{1.15}
\begin{tabular}{>{\raggedright\arraybackslash}p{2.5cm}>{\raggedright\arraybackslash}p{4.0cm}>{\raggedright\arraybackslash}p{6.4cm}}
\toprule
输出类别 & 当前判据 & 解释与验证要求 \\
\midrule
大块异物/超限 & $d_i>50$~mm & 物理阈值明确；需验证尺度与 $d_i$ 误差，适合作为确定性尺寸报警 \\
过小实例/噪声 & $d_i<5$~mm & 当前粗骨料分析的清理口径；是否视为异常取决于具体任务 \\
疑似泥团/非骨料 & $40<d_i\le50$~mm 且 $S_i<0.85$ & 几何启发式候选，只能用于提示复核；需独立泥团/异物真值验证 \\
正常 & 其余未触发规则的实例 & 表示未被当前规则标记，不等价于已证明材质一定为正常骨料 \\
\bottomrule
\end{tabular}
\end{table}

异常标签还会影响最终级配的提交口径。由于第\ref{sec:mass-weighting}节的质量代理会显著放大大粒径实例的权重，确认的超限异物如果直接参与正常骨料级配统计，可能使极少数实例主导 $D_{50}$ 或 $D_{90}$。因此系统同时保留“全部实例结果”和“剔除确认异常后的正常骨料结果”，使评审能够追溯异常处理前后的变化。对于仅由几何启发式得到的疑似泥团，则不应在没有人工或语义证据时无条件删除；正式实验将分别评价异常检测本身的 Precision/Recall，以及异常剔除对最终级配误差的影响。

若比赛数据中泥团或其他非骨料异物占据重要比例，最直接的升级方向是利用实例分割已经生成的颗粒 crop，引入独立的颜色/纹理语义分类器。这样可以保持几何尺寸异常与材质异常两个判断层次解耦：$d_i$ 负责回答“尺寸是否超限”，视觉 crop 分类器负责回答“该实例是否属于骨料”。在当前技术报告中，我们将这一方向明确列为后续增强，而不把尚未验证的语义能力写入现有系统结论。

至此，四项赛题功能在方法层面均已形成明确输出：实例分割提供逐颗粒轮廓，筛分映射提供质量级配，形貌规则提供投影形貌类别，异常模块提供尺寸报警和疑似语义异常候选。下一步的关键不再是继续增加功能，而是建立一套与每项 claim 相匹配的独立证据体系。下一章因此从数据、ground truth、评价指标和消融设计出发，说明如何严格验证这些方法是否真正有效。